\documentclass[11pt,a4paper]{article}

\usepackage[T1]{fontenc}
\usepackage[utf8]{inputenc}
\usepackage{lmodern}
\usepackage{amsmath,amssymb,booktabs,geometry,graphicx,xcolor}
\usepackage{listings}
\usepackage{hyperref}
\geometry{margin=1in}

\definecolor{codegray}{RGB}{245,245,245}
\definecolor{commentgreen}{RGB}{35,120,55}
\lstset{
  language=C,
  basicstyle=\ttfamily\small,
  backgroundcolor=\color{codegray},
  commentstyle=\color{commentgreen},
  keywordstyle=\color{blue},
  numbers=left,
  numberstyle=\tiny\color{gray},
  frame=single,
  breaklines=true,
  columns=fullflexible
}

\title{Grid-Refinement and CFL-Sensitivity Studies of the\\
Quasi-One-Dimensional Compressible Euler Equations\\
in a Convergent--Divergent Nozzle}
\author{Computational Fluid Dynamics Project}
\date{\today}

\begin{document}
\maketitle

\begin{abstract}
This report documents a time-marching, cell-centred finite-volume solution of
the quasi-one-dimensional compressible Euler equations in a convergent--divergent
nozzle. The implementation uses a first-order Rusanov flux and an explicit
CFL-controlled pseudo-time step. The accompanying case files support grid-
refinement studies with 100, 200, and 400 cells and CFL-sensitivity studies with
CFL numbers 0.2, 0.5, 0.8, and 1.0. The solver outputs pressure, Mach number,
flow variables, and residual histories for reproducible analysis.
\end{abstract}

\tableofcontents

\section{Project objective}
The objective is to compute the steady flow through a convergent--divergent
nozzle by advancing the compressible Euler equations in pseudo-time. Two
numerical studies are performed:
\begin{enumerate}
  \item \textbf{Grid refinement:} quantify how the numerical solution changes
        as the number of cells is increased.
  \item \textbf{CFL sensitivity:} compare stability and convergence rate for
        CFL numbers $0.2$, $0.5$, $0.8$, and $1.0$.
\end{enumerate}

The code is intentionally quasi-one-dimensional. This retains the principal
physics of acceleration, choking, and internal shocks while allowing the effect
of numerical resolution and time-step selection to be isolated clearly.

\section{Governing equations}
Let $x$ denote the axial coordinate and $A(x)$ the nozzle cross-sectional area.
The conservative quasi-one-dimensional Euler equations are
\begin{equation}
  \frac{\partial \mathbf U}{\partial t}
  + \frac{\partial \mathbf F}{\partial x}
  = \mathbf S,
\end{equation}
where the area-weighted conservative state is
\begin{equation}
  \mathbf U =
  \begin{bmatrix}
    \rho A\\
    \rho u A\\
    \rho E A
  \end{bmatrix},
\end{equation}
the flux is
\begin{equation}
  \mathbf F =
  \begin{bmatrix}
    \rho u A\\
    (\rho u^2+p)A\\
    u(\rho E+p)A
  \end{bmatrix},
\end{equation}
and the geometric source term is
\begin{equation}
  \mathbf S =
  \begin{bmatrix}
    0\\
    p\,\dfrac{dA}{dx}\\
    0
  \end{bmatrix}.
\end{equation}

For a calorically perfect gas,
\begin{equation}
  E = \frac{p}{\gamma-1}+\frac{1}{2}\rho u^2,
  \qquad
  a = \sqrt{\frac{\gamma p}{\rho}},
  \qquad
  M = \frac{|u|}{a},
\end{equation}
where $a$ is the speed of sound and $M$ is the Mach number.

\section{Nozzle geometry}
The nozzle is one metre long and has its throat at $x_t=0.4$ m. The area is
defined piecewise as a quadratic contraction and expansion:
\begin{equation}
A(x)=
\begin{cases}
A_t+(A_i-A_t)\left(\dfrac{x}{x_t}\right)^2, & 0\leq x\leq x_t,\\[1.2ex]
A_t+(A_e-A_t)\left(\dfrac{x-x_t}{L-x_t}\right)^2,
&x_t < x\leq L.
\end{cases}
\end{equation}

The values used in the study are
\begin{table}[h]
\centering
\caption{Nozzle and gas conditions.}
\begin{tabular}{lll}
\toprule
Quantity & Symbol & Value \\
\midrule
Specific heat ratio & $\gamma$ & $1.4$ \\
Gas constant & $R$ & $287.05$ J/(kg K) \\
Inlet stagnation pressure & $p_0$ & $101325$ Pa \\
Inlet stagnation temperature & $T_0$ & $300$ K \\
Back pressure & $p_b$ & $30000$ Pa \\
Nozzle length & $L$ & $1.0$ m \\
Throat position & $x_t$ & $0.4$ m \\
Inlet area & $A_i$ & $0.020$ m$^2$ \\
Throat area & $A_t$ & $0.010$ m$^2$ \\
Exit area & $A_e$ & $0.016$ m$^2$ \\
\bottomrule
\end{tabular}
\end{table}

\section{Finite-volume discretization}
The domain is divided into $N$ uniform cells with width
\begin{equation}
  \Delta x = \frac{L}{N}.
\end{equation}
Integrating the governing equation over cell $i$ gives
\begin{equation}
  \frac{d\mathbf U_i}{dt}
  = -\frac{\mathbf F_{i+1/2}-\mathbf F_{i-1/2}}{\Delta x}
    + \mathbf S_i.
\end{equation}

The explicit forward-Euler update used by the code is
\begin{equation}
  \mathbf U_i^{n+1}=\mathbf U_i^n
  -\frac{\Delta t}{\Delta x}
   \left(\mathbf F_{i+1/2}^n-\mathbf F_{i-1/2}^n\right)
  +\Delta t\,\mathbf S_i^n.
\end{equation}
The source term is evaluated as
\begin{equation}
  \mathbf S_i =
  \begin{bmatrix}
    0\\p_i A'_i\\0
  \end{bmatrix}.
\end{equation}

\subsection{Rusanov numerical flux}
At each face, the left and right primitive states are used to compute
\begin{equation}
  \mathbf F_{i+1/2} = \frac{1}{2}
  \left(\mathbf F_L+\mathbf F_R\right)
  -\frac{1}{2}s_{\max}\left(\mathbf U_R-\mathbf U_L\right),
\end{equation}
where
\begin{equation}
  s_{\max}=\max\left(|u_L|+a_L,\ |u_R|+a_R\right).
\end{equation}
The Rusanov flux is robust and simple, although it is numerically diffusive,
particularly around shocks.

\subsection{CFL time step}
The time step is selected from the largest characteristic speed in the domain:
\begin{equation}
  \Delta t = \mathrm{CFL}\,
  \min_i\left(\frac{\Delta x}{|u_i|+a_i}\right).
\end{equation}
The CFL number controls the amount of pseudo-time advancement per iteration.
Large values can reduce the number of iterations, but may cause oscillation or
instability for an explicit scheme.

\section{Initial and boundary conditions}
The initial condition is a stationary, uniform flow at the inlet stagnation
state:
\begin{equation}
  \rho_0=\frac{p_0}{RT_0},\qquad u_0=0,\qquad p_0=101325\ \mathrm{Pa}.
\end{equation}

At the inlet, a reservoir state $(\rho_0,0,p_0)$ is supplied to the numerical
flux. At the outlet, density and velocity are extrapolated from the final cell
and pressure is set to the prescribed back pressure $p_b$. These simple boundary
conditions are sufficient for the present study, but characteristic boundary
conditions would be preferable for a production-quality solver.

\section{Algorithm}
The main time-marching loop is:
\begin{enumerate}
  \item Read a case file and construct cell-centre and face areas.
  \item Initialize every cell with the inlet stagnation state.
  \item Recover $(\rho,u,p)$ and calculate the CFL-limited time step.
  \item Construct left and right states at every face.
  \item Evaluate Rusanov fluxes.
  \item Add the area-gradient momentum source term.
  \item Update all conservative variables from step $n$ to step $n+1$.
  \item Record the residual and stop when it is below the requested tolerance.
\end{enumerate}

The residual recorded by the implementation is the average absolute change in
the three conservative variables:
\begin{equation}
  R^n=\frac{1}{N}\sum_{i=1}^{N}\sum_{k=1}^{3}
  \left|U_{i,k}^{n+1}-U_{i,k}^{n}\right|.
\end{equation}

\section{Input cases}
All case files are in the \texttt{cases/} directory. The grid-refinement matrix
holds all physical conditions and CFL number fixed while changing $N$:
\begin{table}[h]
\centering
\caption{Grid-refinement cases.}
\begin{tabular}{llll}
\toprule
Case & Cells & CFL & Output prefix \\
\midrule
\texttt{grid\_100.cfg} & 100 & 0.5 & \texttt{grid\_100} \\
\texttt{grid\_200.cfg} & 200 & 0.5 & \texttt{grid\_200} \\
\texttt{grid\_400.cfg} & 400 & 0.5 & \texttt{grid\_400} \\
\bottomrule
\end{tabular}
\end{table}

The CFL matrix holds the 200-cell grid and all physical conditions fixed:
\begin{table}[h]
\centering
\caption{CFL-sensitivity cases.}
\begin{tabular}{llll}
\toprule
Case & Cells & CFL & Output prefix \\
\midrule
\texttt{cfl\_02.cfg} & 200 & 0.2 & \texttt{cfl\_02} \\
\texttt{cfl\_05.cfg} & 200 & 0.5 & \texttt{cfl\_05} \\
\texttt{cfl\_08.cfg} & 200 & 0.8 & \texttt{cfl\_08} \\
\texttt{cfl\_10.cfg} & 200 & 1.0 & \texttt{cfl\_10} \\
\bottomrule
\end{tabular}
\end{table}

\section{Build and execution}
Compile the solver with:
\begin{lstlisting}[language=bash]
cc -std=c11 -O2 -Wall -Wextra -pedantic nozzle_fvm.c -lm -o nozzle_fvm
\end{lstlisting}

Run one case with:
\begin{lstlisting}[language=bash]
./nozzle_fvm cases/baseline.cfg
\end{lstlisting}

Generate the complete data set with:
\begin{lstlisting}[language=bash]
for f in cases/grid_*.cfg; do ./nozzle_fvm "$f"; done
for f in cases/cfl_*.cfg; do ./nozzle_fvm "$f"; done
\end{lstlisting}

Each case creates two comma-separated files:
\begin{itemize}
  \item \texttt{<prefix>.csv}: $x$, area, density, velocity, pressure, and Mach number;
  \item \texttt{<prefix>\_residual.csv}: iteration number, residual, and time step.
\end{itemize}

\section{Study methodology}
\subsection{Grid refinement}
For each grid, compare the final pressure and Mach profiles. Interpolate the
coarser solution onto the finest-grid coordinates before computing an error:
\begin{equation}
  \|p\|_{L_1} = \frac{1}{N}\sum_i |p_i^{(N)}-p_i^{(\mathrm{fine})}|,
  \qquad
  \|p\|_{L_2} =
  \left[\frac{1}{N}\sum_i
  \left(p_i^{(N)}-p_i^{(\mathrm{fine})}\right)^2\right]^{1/2}.
\end{equation}
The same comparison can be made for Mach number. If an internal shock is
present, estimate its position from the location of the largest magnitude
pressure or Mach-number gradient. Because the present scheme is first-order,
shock position and shock thickness will converge slowly.

\subsection{CFL sensitivity}
For each CFL number, compare:
\begin{itemize}
  \item number of iterations required to reach the residual tolerance;
  \item residual decay as a function of iteration;
  \item minimum pressure and density, as positivity indicators;
  \item final pressure and Mach profiles;
  \item whether the run reaches a stable steady state.
\end{itemize}
The expected trade-off is that larger CFL numbers require fewer iterations when
stable, but may produce oscillations or fail to converge. In the supplied
configuration, CFL $=1.0$ is intentionally included as a stability boundary
test; it does not converge within the configured iteration limit with this
explicit first-order discretization.

\section{Results and analysis}
The figures in this section are generated by \texttt{plot\_results.py} from the
CSV files produced by \texttt{nozzle\_fvm.c}. They are stored in
\texttt{results/visualizations/} so that the report can be rebuilt without
manually editing any image.

\subsection{Grid-refinement results}
All three grid cases converged to the specified residual tolerance. The number
of iterations increases approximately in proportion to the number of cells,
because the explicit time step decreases as $\Delta x$ decreases. The final
iteration counts were 2,656, 5,233, and 10,419 for 100, 200, and 400 cells,
respectively. This is expected for an explicit pseudo-time method: a finer grid
has a smaller stable time step even though it gives a more highly resolved
solution.

\begin{table}[h]
\centering
\caption{Grid-refinement convergence results.}
\begin{tabular}{rrrr}
\toprule
Cells & CFL & Iterations & Final residual \\
\midrule
100 & 0.5 & 2,656 & $9.9952\times10^{-9}$ \\
200 & 0.5 & 5,233 & $9.9602\times10^{-9}$ \\
400 & 0.5 & 10,419 & $9.9858\times10^{-9}$ \\
\bottomrule
\end{tabular}
\end{table}

Figure~\ref{fig:gridprofiles} shows that the pressure and Mach-number profiles
approach a common solution as the grid is refined. The shock, when present, is
spread over several cells by the first-order Rusanov scheme. Therefore, shock
thickness decreases with grid spacing, while the shock location should be
compared using a consistent gradient-based definition.

\begin{figure}[htbp]
  \centering
  \includegraphics[width=0.92\textwidth]{results/visualizations/grid_profiles.png}
  \caption{Pressure and Mach-number profiles for the three grid resolutions.}
  \label{fig:gridprofiles}
\end{figure}

The residual histories are plotted in Figure~\ref{fig:gridresiduals}. The finer
grids require more iterations to reach the same residual threshold, but the
residual histories provide a direct check that all cases converge to a steady
state rather than merely reaching the maximum iteration count.

\begin{figure}[htbp]
  \centering
  \includegraphics[width=0.72\textwidth]{results/visualizations/grid_residuals.png}
  \caption{Residual histories for the grid-refinement cases.}
  \label{fig:gridresiduals}
\end{figure}

\subsection{CFL-sensitivity results}
The CFL study illustrates the stability--efficiency trade-off of explicit time
marching. CFL values 0.2, 0.5, and 0.8 converged. Increasing the CFL number
from 0.2 to 0.8 reduced the iteration count from 12,663 to 3,323. CFL $=1.0$
did not converge within 20,000 iterations and was therefore not used as a
steady-state solution for profile comparisons.

\begin{table}[h]
\centering
\caption{CFL-sensitivity convergence results on the 200-cell grid.}
\begin{tabular}{rrrrl}
\toprule
CFL & Cells & Iterations & Final residual & Status \\
\midrule
0.2 & 200 & 12,663 & $9.9984\times10^{-9}$ & Converged \\
0.5 & 200 & 5,233 & $9.9602\times10^{-9}$ & Converged \\
0.8 & 200 & 3,323 & $9.9586\times10^{-9}$ & Converged \\
1.0 & 200 & 20,000 & $8.47597$ & Not converged \\
\bottomrule
\end{tabular}
\end{table}

Figure~\ref{fig:cflresiduals} shows the faster residual decay obtained with
larger stable CFL values. The CFL $=1.0$ trace is retained because its failure
is an important result: it identifies a stability limitation of the current
explicit, first-order algorithm. It should not be interpreted as a physical
steady solution.

\begin{figure}[htbp]
  \centering
  \includegraphics[width=0.72\textwidth]{results/visualizations/cfl_residuals.png}
  \caption{Residual histories for the CFL-sensitivity study.}
  \label{fig:cflresiduals}
\end{figure}

For the converged CFL cases, the final pressure and Mach profiles are nearly
unchanged, as shown in Figure~\ref{fig:cflprofiles}. This indicates that the
CFL number changes the path to the steady state more strongly than the final
solution, provided the run is stable and sufficiently converged.

\begin{figure}[htbp]
  \centering
  \includegraphics[width=0.92\textwidth]{results/visualizations/cfl_profiles.png}
  \caption{Final pressure and Mach-number profiles for converged CFL cases.}
  \label{fig:cflprofiles}
\end{figure}

\subsection{Baseline flow visualization}
Figure~\ref{fig:baseline} combines the nozzle area, normalized pressure, and
Mach number for the baseline case. It provides a compact numerical check: the
flow accelerates through the convergent section and responds to the divergent
geometry downstream. The present baseline should be interpreted cautiously,
however. Its simple reservoir and fixed-pressure ghost states produce a throat
Mach number above one rather than exactly the ideal choked value $M=1$. This is
a boundary-condition/source-discretization validation issue, not evidence that
the ideal nozzle theory predicts a supersonic throat.

\begin{figure}[htbp]
  \centering
  \includegraphics[width=0.78\textwidth]{results/visualizations/baseline_overview.png}
  \caption{Baseline nozzle geometry and computed flow variables.}
  \label{fig:baseline}
\end{figure}

\subsection{Interpretation}
The results support three conclusions. First, grid refinement changes the
resolution of the flow features and reduces numerical smearing, but it also
increases the computational cost of explicit marching. Second, increasing CFL
from 0.2 to 0.8 improves convergence efficiency for this case without changing
the converged profile materially. Third, CFL $=1.0$ is outside the practical
stable range of this implementation, demonstrating why CFL sensitivity must be
reported alongside the final flow field.

The reported values are numerical results for the specified first-order solver,
not grid-independent physical truth. A higher-order reconstruction, improved
boundary conditions, and a well-balanced source-term treatment would be needed
before making high-accuracy design claims. In particular, before using the
solver for a final choking-position claim, verify that the throat pressure and
Mach number recover the critical relations
\begin{equation}
  \frac{p^*}{p_0}=\left(\frac{2}{\gamma+1}\right)^{\gamma/(\gamma-1)},
  \qquad M^*=1,
\end{equation}
to the expected discretization error.

\section{Code organization}
The implementation is contained in \texttt{nozzle\_fvm.c}:
\begin{description}
  \item[\texttt{nozzle\_area}] evaluates the piecewise nozzle geometry.
  \item[\texttt{area\_slope}] evaluates $dA/dx$ for the geometric source.
  \item[\texttt{primitive\_to\_U}] converts primitive variables to area-weighted conserved variables.
  \item[\texttt{U\_to\_primitive}] recovers density, velocity, and pressure.
  \item[\texttt{rusanov\_flux}] calculates the interface flux.
  \item[\texttt{read\_config}] reads a plain-text case file.
  \item[\texttt{write\_solution}] writes the final flow field to CSV.
  \item[\texttt{main}] performs initialization, time marching, residual output, and termination.
\end{description}

The pressure and density floors in the primitive-variable recovery protect the
calculation against nonphysical intermediate values during unstable runs. They
should be treated as diagnostics, not as a replacement for a more advanced
positivity-preserving scheme.

\section{Limitations and possible extensions}
The solver is intended for a clear numerical-method study rather than a
high-fidelity nozzle design tool. Important limitations are:
\begin{itemize}
  \item first-order spatial accuracy;
  \item explicit forward-Euler time integration;
  \item simple pressure outlet and reservoir inlet conditions;
  \item non-well-balanced treatment of the area source term;
  \item no viscous effects, heat transfer, turbulence, or boundary layers.
\end{itemize}

Natural extensions include MUSCL reconstruction with a limiter, an HLLC flux,
second-order Runge--Kutta time integration, characteristic boundary conditions,
and a well-balanced discretization of the area source term. These extensions
should be added only after the grid and CFL studies are reproducible with the
present baseline solver.

\section{Reproducibility checklist}
\begin{enumerate}
  \item Compile with the command in Section 7.
  \item Run all files in \texttt{cases/}.
  \item Preserve the generated CSV files with the case names.
  \item Record compiler, operating system, grid size, CFL, tolerance, and
        maximum iteration count.
  \item Plot residual histories and final pressure/Mach profiles.
  \item Report non-converged cases separately from converged steady solutions.
\end{enumerate}

\end{document}
